% ------------------------------------------------------------------
%  Capitolo 16 — L'attenzione
%  Parte IV
%  Ricerca di riferimento: ricerca 01 §12; 02 §6.4
%  Voci bibliografiche principali: sana2013, mueller2014, urry2021, delgado2018, beland2016, goodyear2025
%  Epigrafe: ✔ verificata sul testo (vedi ricerca 03 e registro, ricerca 00)
%  Scheda del capitolo: ricerca/06_indice-ragionato.md, capitolo 16
% ------------------------------------------------------------------
\chapter{L'attenzione}
\label{cap:attenzione}

\epigrafe[Non è da nessuna parte chi è dappertutto.]{Nusquam est qui ubique est.}{Seneca, Epistulae ad Lucilium 2, 2}

\iniziale{Q}{uesto capitolo} \segnaposto{apertura: da scrivere — vedi la scheda nell'indice ragionato}.

\section{Una cosa alla volta}

\segnaposto{da scrivere}

\section{Il telefono}

\segnaposto{da scrivere}

\section{Appunti a mano o al computer?}

\segnaposto{da scrivere}

\section{Carta o schermo?}

\segnaposto{da scrivere}

\section{Musica e silenzio}

\segnaposto{da scrivere}

\begin{daricordare}
\segnaposto{il principio del capitolo in due righe}
\end{daricordare}

\begin{domande}
  \item \segnaposto{domanda di recupero su questo capitolo}
  \item \segnaposto{domanda di applicazione a una materia che studi}
  \item \segnaposto{domanda di ripresa su un capitolo precedente}
\end{domande}

\begin{sintesi}
  \item \segnaposto{punto di sintesi}
\end{sintesi}

\finecapitolo
